\documentclass[11pt,a4paper]{article}
\usepackage[margin=1in]{geometry}
\usepackage[T1]{fontenc}
\usepackage{times}
\usepackage{microtype}
\usepackage{enumitem}
\usepackage{booktabs}
\usepackage{tabularx}
\usepackage{amsmath}
\usepackage[dvipsnames]{xcolor}
\usepackage{listings}
\usepackage[hidelinks]{hyperref}

\definecolor{todoblue}{RGB}{0,72,170}
\definecolor{codebg}{RGB}{246,247,249}
\definecolor{donegreen}{RGB}{20,105,55}

\lstset{
  basicstyle=\ttfamily\footnotesize,
  backgroundcolor=\color{codebg},
  breaklines=true,
  frame=leftline,
  framesep=6pt,
  rulecolor=\color{todoblue!40},
  xleftmargin=10pt,
  columns=fullflexible,
  keepspaces=true,
}

\setlength{\parindent}{0pt}
\setlength{\parskip}{0.5em}

\title{\textbf{H-ZKA: Status and Outstanding Work}\\[0.3em]
\large What the third revision corrected, and what remains}
\author{Working notes for the author team}
\date{22 August 2026}

\begin{document}
\maketitle

\section*{Where things stand}

The Associate Editor's report identified five internal-validity issues and
judged the previous TODO list ``valuable but misprioritized for acceptance.''
That judgement was correct. This round therefore did the five corrections and
deferred the extensions. This document records both.

\vspace{0.4em}
\begin{tabularx}{\textwidth}{@{}p{0.5cm}>{\raggedright\arraybackslash}p{5.0cm}p{2.1cm}>{\raggedright\arraybackslash}X@{}}
\toprule
& \textbf{Item} & \textbf{Status} & \textbf{Outcome} \\
\midrule
1 & Two-stage prover accounting & \textcolor{donegreen}{\textbf{Done}} &
$+17\%$ workers at $k=100$, not $-83\%$. Three claims withdrawn \\
2 & Constant-size public-input interface & \textcolor{donegreen}{\textbf{Done}} &
$O(\sqrt k)$ restored in every term; exact gas model replaces a mocked
measurement \\
3 & Per-cluster BFT conditioning & \textcolor{donegreen}{\textbf{Done}} &
Captured clusters now modelled; $50\%$ claim withdrawn \\
4 & Matched leakage baseline & \textcolor{donegreen}{\textbf{Done}} &
$3.0\times$ worst case, not $56.9\times$; benefit traced to the interface \\
5 & Failover deadline semantics & \textcolor{donegreen}{\textbf{Done}} &
Lost-round rate now reported: $1.3\%$ at a $10\%$ crash rate \\
\midrule
6 & Execute the 20M circuit, matched toolchain & Outstanding &
Removes the last predicted headline number \\
7 & Trace-calibrated simulation inputs & Outstanding &
Improves external validity; does \emph{not} change the evidence class \\
8 & Alternative proving backend & Outstanding &
Optional; the AE noted Reviewer 3 asked for this \emph{or} larger validation \\
9 & Portal metadata and release tag & Author only & Required before upload \\
\bottomrule
\end{tabularx}

\section{Execute the 20-million-constraint circuit, with a matched toolchain}
\label{sec:todo-20m}

\subsection*{Why the previous plan was not sufficient}

The earlier version of this document proposed a Circom and snarkjs run. The
Associate Editor pointed out that the manuscript describes an xjsnark
implementation and a separate Rust/Arkworks profiling harness, so a run in a
third toolchain would be read as a \emph{new implementation} rather than as
validation of the reported model. That objection is correct and it changes the
requirement: the run must match the manuscript, or the manuscript must be
changed to match the run.

\subsection*{What must match}

Five things, and all five should be recorded in the artifact before the run
starts:

\begin{enumerate}[leftmargin=*,itemsep=3pt]
  \item \textbf{Statement.} Verify $B_{\max}=15$ inner Groth16 proofs and check
  $\mathsf{StateTransition}(\mathsf{rt}^{\mathrm{old}}_j,\mathsf{tx}_j)=
  \mathsf{rt}^{\mathrm{new}}_j$ for each slot.
  \item \textbf{Public-input interface.} Exactly three public inputs: the
  cluster commitment, the cluster id, and the round. This is the interface
  Algorithm~1 now specifies, and the run must use it, not the per-chain variant.
  The commitment binding, $B_{\max}-1$ Poseidon compressions, must be inside
  the circuit and counted in the constraint total.
  \item \textbf{Security level.} 128-bit, BN254, matching the manuscript.
  \item \textbf{Toolchain.} Rust/Arkworks with the \texttt{parallel} feature,
  which is what produced the nine measured points in Table~19. Using it keeps
  the new datum on the same regression line rather than starting a new one.
  \item \textbf{Host.} The AsusL40 node (Xeon Gold 6538Y+, 16 vCPU, 200 GiB),
  the same host as the nine measured points. Provision at least 128 GiB free:
  the measured profile is 3 GiB per million constraints, so 20M needs about
  60 GiB for the prover alone plus proving key and witness.
\end{enumerate}

\subsection*{Sequence}

\begin{lstlisting}[caption={Run from the Arkworks harness, not from a fresh Circom build}]
# 1. Build the aggregation circuit at B_max = 15 with the commitment interface.
#    Record the constraint count and the source hash.
cargo run --release --features parallel --bin build_agg -- \
    --slots 15 --public-inputs commitment --out build_agg/

# 2. Record identity before proving.
sha256sum build_agg/agg.r1cs build_agg/agg_pk.bin build_agg/agg_vk.bin
cargo tree --locked > build_agg/deps.lock.txt

# 3. Witness generation, timed separately from proving.
/usr/bin/time -v cargo run --release --features parallel --bin witness -- \
    --circuit build_agg/agg.r1cs --input input_full.json \
    --out witness.bin

# 4. Prove, three independent runs, peak RSS captured each time.
for i in 1 2 3; do
  /usr/bin/time -v cargo run --release --features parallel --bin prove -- \
      --pk build_agg/agg_pk.bin --witness witness.bin \
      --out proof_$i.bin 2> prove_$i.time
done

# 5. Verify, and record the public inputs actually accepted.
cargo run --release --bin verify -- \
    --vk build_agg/agg_vk.bin --proof proof_1.bin --public public_1.json
\end{lstlisting}

Store everything in \texttt{HZKA/result/circuit\_20M/} and add a manifest row.

\subsection*{What to change in the manuscript afterwards}

\begin{enumerate}[leftmargin=*,itemsep=3pt]
  \item \textbf{Table 8, p.~18.} Change the 20M row from \emph{Predicted} to
  \emph{Measured}.
  \item \textbf{Table 19.} Replace the interval row with the measured mean and
  CI over three runs.
  \item \textbf{Section 7.9, Table 21.} This becomes a validation of the three
  model forms rather than a risk bound. Add a column giving each model's error
  against the measured value. Keep the section: knowing which extrapolation
  form was right is more useful than deleting the analysis.
  \item \textbf{Section 7.5.1, Tables 14--15.} Recompute worker counts, memory,
  critical path, and the latency-parity $k$ from Eq.~(42) with the measured
  $S_{\mathrm{agg}}$. If the measured value is below about $75$\,s, parity moves
  inside $k\leq200$ and the latency claim can be revisited; the manuscript
  currently claims none.
  \item \textbf{Table 42 and Section 8.2.} Recompute the $+17\%$ and $+29\%$
  overheads.
  \item \textbf{Abstract, Section 9 item (1), Section 10.} Remove the
  prediction language.
\end{enumerate}

If the measured value is \emph{worse} than predicted, the corrected claims
still hold in direction; only the magnitudes move. Report it either way.

\section{Trace-calibrated simulation inputs}
\label{sec:todo-traces}

\subsection*{The labelling point, accepted}

The Associate Editor noted that feeding measured RTT, flow correlation, and
churn distributions into the simulator produces \textbf{trace-calibrated
simulation}, not measured protocol behavior, and that the evidence class cannot
change unless the protocol itself is executed on those traces. We accept this
without qualification, and the manuscript's evidence vocabulary now uses that
term. Plan the work accordingly: it improves external validity and does not
promote anything from \emph{simulated} to \emph{measured}.

\subsection*{What to do}

The simulator already accepts all three inputs; only a loader is missing.

\begin{enumerate}[leftmargin=*,itemsep=3pt]
  \item \textbf{Real RTT matrix.} Pairwise RTT between public RPC endpoints, or
  between cloud regions. Feed as \texttt{Topology.rtt\_ms}.
  \item \textbf{Real flow correlation.} Derive $\rho_{ab}$ from observed bridge
  volume between chain pairs. Feed as \texttt{Topology.corr}. This matters most
  for Section~7.14, because the whole $\eta$ trade-off depends on how far
  geography and flow diverge in practice.
  \item \textbf{Real churn.} Fit the offline/online process to observed
  validator uptime rather than the geometric process now used.
\end{enumerate}

Add \texttt{--topology-file} to each driver and keep the synthetic path as the
default so published numbers stay reproducible. Add a fourth evidence status,
\emph{trace-calibrated}, to Table~8 rather than reusing \emph{measured}.

\subsection*{Second priority: extending protection past the BFT bound}

Section~7.12.1 now shows that the per-cluster Byzantine condition, not the
global ratio, is the binding constraint, and that a single captured cluster
blocks global round completion. This is the most consequential open problem for
deployment, and three directions are worth testing in the simulator before any
of them reaches the paper:

\begin{enumerate}[leftmargin=*,itemsep=3pt]
  \item \textbf{Larger balanced clusters.} Table~35 shows capture exposure is
  not monotone in $B$ because $\lceil B/3\rceil/B$ is a step function. Sweep
  $B$ jointly with a balance constraint on the $k$-medoid partition and find
  the configuration that maximizes the safety margin at acceptable circuit size.
  \item \textbf{Shorter epochs.} $E_{\mathrm{epoch}}=100$ lets a capture persist
  for up to an epoch. Sweep $E_{\mathrm{epoch}}$ against reassignment churn,
  which Table~30 already measures.
  \item \textbf{Partial-round completeness.} The current rule makes a round
  incomplete if any slot is missing. A rule that commits the fresh slots and
  carries the missing ones forward would degrade far more gracefully. This is a
  protocol change, not a parameter change, and needs its own safety argument.
\end{enumerate}

\section{Alternative proving backend}
\label{sec:todo-backend}

Still worth doing, still not required for acceptance: the Associate Editor
noted that Reviewer~3 asked for larger validation \emph{or} another backend,
and item~\ref{sec:todo-20m} satisfies the first.

If undertaken, target \textbf{Nova} or HyperNova folding rather than Halo2.
Folding amortizes $B_{\max}$ inner verifications into repeated applications of
one step circuit, which matches Algorithm~1 exactly, and it removes the
per-capacity trusted setup that Section~8.4 identifies as the largest
planned-maintenance event in the architecture. Hold fixed the same five things
listed in Section~\ref{sec:todo-20m}, and report EVM verification gas as well:
a folding backend usually trades cheaper proving for a more expensive on-chain
verifier, and since the paper's claim is now explicitly about on-chain cost,
that would matter.

\section{Author-only actions before upload}
\label{sec:todo-portal}

\begin{enumerate}[leftmargin=*,itemsep=3pt]
  \item \textbf{Tagged release.} The Associate Editor asked for an immutable
  release for the reproducibility spot-check. The corrected code is committed as
  \texttt{a103cad9}. Tag and push it:
  \begin{lstlisting}
git tag -a v3.0-jsa-r3 -m "JSA third revision" a103cad9
git push origin v3.0-jsa-r3
  \end{lstlisting}
  Then add the tag to the \texttt{hzka\_artifact} bibliography entry so the
  cited URL resolves to this exact state.
  \item \textbf{ORCID identifiers} for all four authors if the portal requests
  them.
  \item \textbf{Prior-conference disclosure.} If asked, supply the exact
  conference title, venue, year, and DOI. These are not in the project files
  and were deliberately not invented.
  \item \textbf{Declarations.} Reconfirm author approval, exclusivity, and the
  no-conflict declaration. Funding is already stated (NCM2025-26-01).
\end{enumerate}

\section{Smaller items}

\begin{enumerate}[leftmargin=*,itemsep=4pt]
  \item \textbf{Reputation-DoS equilibrium.} Section~9 still declines to claim a
  Nash equilibrium for frivolous-challenge deterrence. The simulator now has
  what is needed: add a challenger agent with a cost model and sweep the
  challenge bond against the accuser's expected gain.

  \item \textbf{Presence-bitmap leakage.} Section~7.15 excludes the bitmap
  because it publishes participation rather than state change. Quantifying what
  an observer learns from participation alone would complete the picture;
  \texttt{exp\_leakage.py} extends to it by changing the secret and adding the
  bitmap to the observation vector.

  \item \textbf{Float placement.} Two or three wide tables may still sit two
  pages from their sections in the final typeset version. This is an
  \texttt{elsarticle} two-column constraint; if the production editor objects,
  split the widest table by group.

  \item \textbf{Contract alignment.} \texttt{AuditContractV2.acceptClusterCommit}
  still takes \texttt{chainIds[]} and \texttt{newRoots[]} with an $O(B^2)$
  membership loop, which is the per-chain interface the manuscript no longer
  specifies. Aligning the contract with Algorithm~1's commitment interface is
  not needed for the paper's claims, since those now rest on the gas model, but
  it should be done before anyone treats the contract as a reference
  implementation.
\end{enumerate}

\section*{Where everything lives}

\begin{tabularx}{\textwidth}{@{}>{\raggedright\arraybackslash}p{6.6cm}>{\raggedright\arraybackslash}X@{}}
\toprule
\textbf{Path} & \textbf{Contents} \\
\midrule
\texttt{main.tex} / \texttt{main.pdf} & Clean manuscript, 38 pages \\
\texttt{main\_marked\_up.tex} / \texttt{.pdf} & Marked copy, generated, never hand-edited \\
\texttt{make\_marked.py} & Generator for the marked copy \\
\texttt{response.tex} / \texttt{.pdf} & Response to the Associate Editor, 8 pages \\
\texttt{coverletter.tex} / \texttt{.pdf} & Cover letter, one A4 page \\
\texttt{report\_audit.md} & Human-maintained readiness audit \\
\texttt{.revision\_backup\_} \texttt{20260822\_r3/} & The version the Associate Editor reviewed \\
\midrule
\multicolumn{2}{@{}l}{\textit{In the public artifact repository, commit \texttt{a103cad9}:}} \\
\texttt{HZKA/scripts/revision2/} & Simulator and ten experiment drivers \\
\texttt{HZKA/result/revision2/} & Per-seed outputs for every derived table \\
\bottomrule
\end{tabularx}

\vspace{0.8em}

To regenerate everything:

\begin{lstlisting}
cd HZKA/scripts/revision2 && bash run_all.sh
\end{lstlisting}

Base seed \texttt{20260822}, 30 seeds per cell. Re-running reproduces every
published number exactly.

\end{document}
